\paragraph{E.3o (Known, with the integral consequence proved: the Walker
kernel is Saito's product relation).}
We work over $\mathbb C$ with integral Chow groups. For a smooth
projective variety $V$, let
\[
 \mathcal A^r(V)=\mathrm{CH}^r(V)_{\mathrm{alg}},\qquad
 \mathcal H^r(V)=\ker\bigl(\mathrm{CH}^r(V)
                 \longrightarrow H^{2r}(V,\mathbb Z(r))\bigr).
\]
For two adequate relations $\mathcal E,\mathcal F$, we use
Saito's product: $(\mathcal E*\mathcal F)^i(X)$ is the subgroup
generated by
\[
 (p_X)_*(a\cdot b),\qquad
 a\in\mathcal E^r(T\times X),\quad
 b\in\mathcal F^s(T\times X),\quad
 r+s=\dim T+i,
\]
where $T$ is smooth projective and $p_X:T\times X\to X$.
This is an adequate relation by
\citet[Section 1.3 and Lemma 1.3.1, pp.~293--294]{cySaito1992}.
For an adequate relation $\mathcal E$, we also set
\[
 \langle\mathcal E\rangle_0^i(X)
 =\sum_{T,\,W\in\mathrm{CH}^i(T\times X)}
       W_*\bigl(\mathcal E^{\dim T}(T)\bigr).
\]
This is exactly the zero-cycle-generated relation of
\citet[Section 1.7 and Lemma 1.7.1, p.~296]{cySaito1992}.
The following equalities are integral:
\begin{equation}\label{eq:walker-saito-product}
 \ker\psi_W^i(X)=(\mathcal A*\mathcal H)^i(X),
 \qquad
 N^i(X)=\langle\mathcal H^{*2}\rangle_0^i(X).
\end{equation}
Consequently the general complex converse in E.3 is equivalent
to the precise generation assertion
\begin{equation}\label{eq:remaining-zero-generation}
 \boxed{\quad
 \langle\mathcal H^{*2}\rangle_0^i(X)
       =(\mathcal A*\mathcal H)^i(X).
 \quad}
\end{equation}
E.3u below refutes \eqref{eq:remaining-zero-generation}
in general, by an explicit codimension-four family.

\begin{proof}
We first identify exactly the statement we take from Saito.
His Example 1.8, equation (1.8.2) on p.~296, defines
\[
 \operatorname{Gr}^{l}\mathrm{CH}(V)
 =\frac{\langle\mathcal H^{*l}\rangle_0\mathrm{CH}(V)}
 {\bigl(\langle\mathcal H^{*l}\rangle_0
                 \cap\mathcal H^{*(l+1)}\bigr)\mathrm{CH}(V)}.
\]
In particular its first step is
$\mathcal A/(\mathcal A\cap\mathcal H^{*2})$, since
$\langle\mathcal H\rangle_0=\mathcal A$ by his Example 1.9,
p.~297. The denominator is an intersection; it is not
$\langle\mathcal H^{*(l+1)}\rangle_0$.
Thus Theorem 6.2 does not by itself identify the manuscript's
quotient $P^i(X)$ with an abelian variety.

The needed positive statement is explicit in
\citet[proof of Theorem 6.2, p.~323; Remark 6.4, p.~325]{cySaito1992}:
\[
 K_{\mathrm{AJ}}^i(X)/(\mathcal A*\mathcal H)^i(X)
 \quad\text{is finite},\qquad
 K_{\mathrm{AJ}}^i(X)=
       \ker\bigl(\mathcal A^i(X)\to J_a^i(X)(\mathbb C)\bigr).
\]
For clarity, Saito introduces
$'\operatorname{Gr}^{1}\mathrm{CH}^i(X)
=\mathcal A^i(X)/(\mathcal A*\mathcal H)^i(X)$ in
Lemma 6.2.1, p.~321. Its map to $J_a^i(X)$ is covered by
an abelian variety whose cohomological map is an isomorphism
over $\mathbb Q$; the composition is an isogeny. The
surjectivity proved on pp.~322--323 makes its kernel finite.
This is an assertion about the integral groups just displayed,
not a statement only after tensoring by $\mathbb Q$.

We record why the two divisibility steps have the required
scope. The group $(\mathcal A*\mathcal H)^i(X)$ is divisible:
in each generator $(p_X)_*(a b)$ we divide the algebraically
trivial factor $a$ by the chosen positive integer, leaving
$b$ fixed. This also proves directly the instance of Saito's
Lemma 1.10, p.~297, that we use. Ordinary Abel--Jacobi
vanishes on this group. Equivalently this follows from the
inclusion in Remark 6.4, or by representing $a$ by a curve
family and using that multiplication by $b$ makes the
correspondence's cohomological action zero.

Write
\[
 \mathcal A^i(X)\xrightarrow{\psi_W^i}J_W^i(X)(\mathbb C)
     \xrightarrow{\alpha}J_a^i(X)(\mathbb C),
\]
where $\alpha$ is the Walker isogeny. The image of
$(\mathcal A*\mathcal H)^i(X)$ under $\psi_W^i$ lies in
the finite group $\ker\alpha(\mathbb C)$. A homomorphism
from a divisible group to a finite group is zero; hence
$(\mathcal A*\mathcal H)^i(X)\subset\ker\psi_W^i$.
The quotient
\[
 \ker\psi_W^i/(\mathcal A*\mathcal H)^i(X)
\]
is therefore a subgroup of Saito's finite quotient. It is
also divisible, because the whole Walker kernel is divisible
by \citet[Theorem 2.3, author p.~12]{cyACMVWalker2022}.
It is zero. This proves the first equality in
\eqref{eq:walker-saito-product}.

For a smooth projective $t$-fold $T$, Saito's Corollary 6.3
and Remark 6.4, p.~325, specialize to the integral equality
\[
 (\mathcal A*\mathcal H)^t(T)
 =\ker\bigl(\mathrm{CH}_0(T)^0
                      \to\operatorname{Alb}(T)(\mathbb C)\bigr).
\]
Here the absence of a finite discrepancy uses the injectivity
of the zero-cycle Albanese map on torsion, explicitly used
in his proof on p.~323. Moreover
$\mathcal A*\mathcal H\subset\mathcal H^{*2}$, and
$\mathcal H^{*2}$ is Abel--Jacobi trivial by Example 1.11,
pp.~297--298. On zero-cycles, ordinary Abel--Jacobi is the
Albanese map. These two inclusions give
\[
 (\mathcal H^{*2})^t(T)
 =\ker\bigl(\mathrm{CH}_0(T)^0
                      \to\operatorname{Alb}(T)(\mathbb C)\bigr).
\]
Taking the correspondence images in the definition of
$\langle\mathcal H^{*2}\rangle_0$ is now precisely the
definition of $N^i(X)$. This proves the second equality.
Together with E.3m--n the two equalities prove the asserted
fixed-target equivalence.
\end{proof}

The same product $\mathcal A*\mathcal H$ is obtained if
the second factor is required only to have zero rational
Betti class. Indeed an integral cycle $b$ with zero rational
class has $mb\in\mathcal H$ for some $m>0$; writing
$a=ma'$ within the divisible group $\mathcal A$ gives
$ab=a'(mb)$. This observation does not replace the
integral homological relation in the displayed definition
of $\langle\mathcal H^{*2}\rangle_0$.

\paragraph{E.3p (Proved: only surfaces are needed in the
generation test).}
For every smooth projective complex $X$ and every $i$,
\begin{equation}\label{eq:surface-generated-N}
 N^i(X)=
 \sum_{\substack{S/\mathbb C\ \mathrm{smooth\ projective\ surface}\\
                 W\in\mathrm{CH}^i(S\times X)}}
 W_*\ker\bigl(\mathrm{CH}_0(S)^0
                   \to\operatorname{Alb}(S)(\mathbb C)\bigr)
 \quad\text{integrally}.
\end{equation}
In particular \eqref{eq:remaining-zero-generation} is the
assertion that every sum of the products in E.3o is a finite
sum of these surface-kernel images. This is a restriction
on the dimension of the generating varieties, not an
assertion that their Albanese kernels are themselves known.
It agrees with
\citet[Lemma 7.1.1, pp.~326--327]{cySaito1992}; we give
the direct argument appropriate to the present definition.

\begin{proof}
Consider one generator $W_*\nu$ of $N^i(X)$, where
$T$ is smooth projective connected and $\nu\in
\mathrm{CH}_0(T)^0$ has zero Albanese class. Choose an
integral zero-cycle representative of $\nu$ and let $F$
be its finite support. If $\dim T\ge3$, sufficiently
positive successive hyperplane sections through $F$
give a smooth connected complete-intersection surface
$j:S\hookrightarrow T$. To see that the imposed points
cause no obstruction, take powers of a very ample bundle
whose sections vanishing on $F$ generate the conormal
directions at every point of $F$; a general section is
then smooth there, and Bertini applies away from $F$.
Repeat until the dimension is two.

Integral weak Lefschetz gives
$H_1(S,\mathbb Z)\simeq H_1(T,\mathbb Z)$, compatibly
with the Hodge structures and their periods. Thus
\[
 j_*:\operatorname{Alb}(S)\xrightarrow{\sim}
                         \operatorname{Alb}(T)
\]
is an isomorphism of complex abelian varieties, rather
than merely an isogeny. The same formal sum of the points
of $F$ defines a degree-zero cycle $\nu_S$ on $S$.
Its Albanese class maps to that of $\nu$, so is zero.
The refined pullback $W_S=(j\times1_X)^*W$ is an
integral class in $\mathrm{CH}^i(S\times X)$, and
the compatibility of correspondences with proper
pushforward gives
\[
                   W_*\nu=(W_S)_*\nu_S.
\]
No transversality of a chosen representative of $W$
is required for this equality of Chow operations.

If $\dim T=2$ there is nothing to do. In dimension one
the Albanese kernel is zero by Abel's theorem, and
in dimension zero the connected degree-zero Chow group
is zero. We have expressed every generator using a
surface; the reverse inclusion is immediate from the
definition of $N$.
\end{proof}

The literal ordinary Abel--Jacobi equality is still
false in the Ottem--Suzuki example of E.3h--k. The
zero-cycle generation assertion in
\eqref{eq:remaining-zero-generation}, equivalently the
surface generation in \eqref{eq:surface-generated-N}
of all $\mathcal A*\mathcal H$ classes, is refuted by E.3u.
A universal
regular-homomorphism theorem does not assert this
equality of kernels: Samuel's abelian equivalence is
defined by the intersection of all regular kernels,
whereas $N$ is generated by the specified correspondences.
The distinction is explicit in
\citet[Section 1.1, pp.~7--8]{cyBarbieriViale2001}
and \citet[Section 4.1, pp.~13--16]{cyHernandezPelaez2025}.
\paragraph{E.3q (Proved under stated hypotheses: the codimension-two
case with surface-supported zero-cycles).}
Let $X/\mathbb C$ be smooth projective of dimension $n\ge2$,
and suppose that $\mathrm{CH}_0(X)$ is supported on a closed
subset of dimension at most two, in the sense used in E.3l.
Then
\[
 P^2(X)\ \xrightarrow{\ \sim\ }\ J_W^2(X)(\mathbb C)
          \ \xrightarrow{\ \sim\ }\ J_a^2(X)(\mathbb C)
\]
integrally. Thus the family converse and the integral equality
of ordinary Abel--Jacobi and Albanese equivalence hold in
codimension two under this hypothesis. Together with E.3l this
settles both codimensions $2$ and $n-1$ under the same support
assumption; for $n=3$ they are the same codimension.

\begin{proof}
We transpose the actual decomposition of the diagonal
constructed in E.3l. After resolving its support and clearing
one integer, its action on $\mathrm{CH}^2(X)$ has the form
\begin{equation}\label{eq:transpose-surface-factorization}
 M\,\mathrm{id}
   =\sum_j (G_j)_* f_j^*
          +\sum_k (g_k)_* (H_k)_* .
\end{equation}
Here $M>0$, the smooth projective $S_j$ have dimension
$s_j\le2$, and the smooth projective $D_k$ have dimension
$n-1$. The proper maps and correspondences have the types
\[
 \begin{gathered}
 f_j:S_j\longrightarrow X,\qquad
 G_j\in\mathrm{CH}^{s_j}(S_j\times X),\\
 g_k:D_k\longrightarrow X,\qquad
 H_k\in\mathrm{CH}^{n-1}(X\times D_k).
 \end{gathered}
\]
Consequently the first summand in
\eqref{eq:transpose-surface-factorization} acts through
$\mathrm{CH}^2(S_j)$, and the second acts through
$\mathrm{CH}^1(D_k)$:
\[
 \begin{gathered}
 \mathrm{CH}^2(X)\xrightarrow{f_j^*}
       \mathrm{CH}^2(S_j)\xrightarrow{(G_j)_*}\mathrm{CH}^2(X),\\
 \mathrm{CH}^2(X)\xrightarrow{(H_k)_*}
       \mathrm{CH}^1(D_k)\xrightarrow{(g_k)_*}\mathrm{CH}^2(X).
 \end{gathered}
\]
All these operations are integral and preserve algebraic
equivalence and $N$. The same support-resolution argument
as in E.3l supplies the integer $M$; transposition does not
introduce any additional denominator.

Let $z\in\mathcal A^2(X)$ have zero ordinary Abel--Jacobi
class. Correspondence functoriality implies that
$f_j^*z$ has zero Albanese class when $s_j=2$.
It is zero for dimension reasons when $s_j<2$.
Likewise $(H_k)_*z\in\operatorname{Pic}^0(D_k)$ has
zero divisor Abel--Jacobi class, so is zero as a
divisor class. Formula \eqref{eq:transpose-surface-factorization}
therefore gives $Mz\in N^2(X)$. It follows that
$\ker\psi_W^2/N^2(X)$ is killed by $M$.
It is a divisible group by ACMV's Theorem 2.3, and hence
zero. This proves the first isomorphism.

For the second, Saito's Corollary 6.3 and Remark 6.4,
p.~325, give
$\ker\psi_{\mathrm{AJ}}^2=(\mathcal A*\mathcal H)^2(X)$
integrally; E.3o identifies this with $\ker\psi_W^2$.
Since $\psi_W^2$ is surjective, the Walker isogeny has
zero kernel and is an isomorphism. This argument uses
the exact codimension-two torsion theorem incorporated
in Saito's Corollary 6.3; it does not extend the same
ordinary-Jacobian conclusion to codimension $n-1$ in
higher dimensions.
\end{proof}


